\chapter{图论进阶：并查集、最短路和状态图}

图论进阶题的难点通常不在 BFS 或 DFS 的基本写法，而在选择正确的图模型。无向连通性适合并查集；无权最短路适合 BFS；边权非负的最短路适合 Dijkstra；存在状态压缩或多维状态时，要先定义状态节点，再决定边和代价。很多题目表面不是图，例如密码锁、单词变化、网格体力消耗、课程安排、账户合并，只要存在“状态”和“一步转移”，就可以用图论语言重写。

本章先从并查集讲动态连通性，再讲拓扑排序、Dijkstra、0-1 BFS、网格代价、强连通分量、桥、二分图匹配、网络流和状态图。第 8 章已经负责树图基础、普通 BFS/DFS 和基础建图，本章不重复那些入口知识，而是专门分析“什么时候普通 BFS 不够”“什么时候只问连通性不用最短路”“什么时候状态必须包含钥匙、步数、剩余资源或访问集合”“什么时候图结构本身需要被压缩成分量或残量网络”。每个案例都按暴力搜索、状态压缩、剪掉重复、证明边界的顺序展开。

\section{并查集、连通性和动态合并}

并查集解决的是动态连通性。它不回答路径是什么，也不回答最短距离是多少；它只维护“两个元素当前是否属于同一集合”。暴力做法在每次合并后用 DFS 或 BFS 检查连通性，会重复扫描大量边。并查集把每个集合表示成一棵树，用根节点作为代表。路径压缩让查询时把沿途节点直接挂到根上，按大小合并让小树挂到大树下，两者结合后，均摊复杂度接近常数。

\begin{lstlisting}[language=C++,caption={并查集：迭代 find、路径压缩和按大小合并}]
class DisjointSetUnion {
public:
    explicit DisjointSetUnion(int n)
        : parent_(n), size_(n, 1)
    {
        std::iota(this->parent_.begin(), this->parent_.end(), 0);
    }

    int find(int x)
    {
        int root = x;
        while (this->parent_[root] != root) {
            root = this->parent_[root];
        }
        while (this->parent_[x] != x) {
            const int next = this->parent_[x];
            this->parent_[x] = root;
            x = next;
        }
        return root;
    }

    bool unite(int a, int b)
    {
        int root_a = this->find(a);
        int root_b = this->find(b);
        if (root_a == root_b) {
            return false;
        }
        if (this->size_[root_a] < this->size_[root_b]) {
            std::swap(root_a, root_b);
        }
        this->parent_[root_b] = root_a;
        this->size_[root_a] += this->size_[root_b];
        return true;
    }

private:
    std::vector<int> parent_;
    std::vector<int> size_;
};
\end{lstlisting}

这里的 \code{find} 是迭代写法，避免递归深度问题。第一次循环找根，第二次循环压缩路径。 \code{unite} 返回 \code{false} 表示两个元素本来已经连通，合并会形成环。这个返回值在很多题里正好就是答案判断条件。成员访问使用 \code{this->}，让读者清楚哪些状态属于并查集对象。

冗余连接是并查集的标准例题。题目给一棵树额外加一条边，要求找出导致环的边。暴力方法是每加入一条边，就在已有图中检查这两个端点是否已经连通；如果已经连通，再加边就会成环。并查集正好维护这个关系：扫描边时，如果 \code{unite(u, v)} 返回 \code{false}，说明这条边连接了同一集合内的两个点，它就是冗余边。

\begin{lstlisting}[language=C++,caption={LeetCode 684 冗余连接：合并失败的边就是成环边}]
std::vector<int> find_redundant_connection(
    const std::vector<std::vector<int>>& edges)
{
    DisjointSetUnion dsu(static_cast<int>(edges.size()) + 1);
    for (const std::vector<int>& edge : edges) {
        const int u = edge[0];
        const int v = edge[1];
        if (!dsu.unite(u, v)) {
            return edge;
        }
    }
    return {};
}
\end{lstlisting}

这题要注意节点编号通常从 1 开始，所以并查集大小开到 \code{edges.size() + 1}。如果题目节点编号不连续，就需要哈希映射到连续编号。复杂度近似 \complexity{O(n)}。面试表达中要把树的性质说出来：一棵 n 个节点的树有 n-1 条边且无环；额外加入一条边后，第一次发现两个端点已经连通的边，就是让无环结构变成有环结构的边。

省份数量也是连通块问题。输入是邻接矩阵，\code{isConnected[i][j] == 1} 表示城市 i 和 j 直接连接。DFS/BFS 可以做，并查集也可以做。暴力思路可能是从每个城市出发搜索所有可达城市并统计，但会重复访问。并查集做法扫描矩阵上三角，把直接连接的城市合并，最后统计根节点个数。

\begin{lstlisting}[language=C++,caption={LeetCode 547 省份数量：邻接矩阵合并连通城市}]
int find_circle_num(const std::vector<std::vector<int>>& is_connected)
{
    const int n = static_cast<int>(is_connected.size());
    DisjointSetUnion dsu(n);

    for (int i = 0; i < n; ++i) {
        for (int j = i + 1; j < n; ++j) {
            if (is_connected[i][j] == 1) {
                dsu.unite(i, j);
            }
        }
    }

    int provinces = 0;
    for (int city = 0; city < n; ++city) {
        if (dsu.find(city) == city) {
            ++provinces;
        }
    }
    return provinces;
}
\end{lstlisting}

邻接矩阵是对称的，因此只扫描上三角即可。这里并查集的优势不是绝对性能，而是模型清晰：省份就是连通分量，直接连接就是合并操作。若题目还要求列出每个省份包含哪些城市，则并查集统计完根后还要按根分组；若题目要求最短路径，并查集就不够，因为它丢掉了路径长度信息。

\section{最短路算法族：从等权边到非负边和负边}

Dijkstra 处理非负边权最短路。无权图 BFS 成立，是因为每条边代价相同，按层访问保证第一次到达最短；带权图中，一条边可能代价很大，普通 BFS 不再可靠。Dijkstra 的贪心点是：每次从尚未确定的节点中取当前距离最小者，它的最短距离已经确定。这个结论依赖所有边权非负，因为后续绕路不可能通过负边把距离降得更低。

最短路题要先区分三件事。第一，路径代价如何合成：普通最短路把边权相加，最小体力路径取路径上最大边权，最大概率路径把概率相乘，迷宫滚动题把一次滚动距离累加。第二，扩展顺序由什么保证正确：等权图用队列的层次，非负加法代价用小顶堆，0/1 权用双端队列，限制边数用分层松弛。第三，一个状态什么时候可以永久确定：BFS 中第一次入队通常就是最短层；Dijkstra 中从堆里弹出且不是过期距离时才确定；Bellman-Ford 中没有“弹出即确定”，只有经过固定轮数后的阶段性最优。

从暴力角度看，最短路最直接的想法是枚举所有从起点到终点的路径，再取代价最小者。这个想法很快爆炸，因为有环图中路径数量可以无限，哪怕限制为简单路径也可能指数级。BFS、Dijkstra、Bellman-Ford 的共同目标，都是避免重复枚举路径，而是维护“到某个状态目前已知的最好代价”。不同算法的差别在于：它们用什么规则保证某些状态不用再被未来路径改进。

\begin{center}
\begin{tabular}{llll}
\toprule
边权和约束 & 推荐模型 & 扩展顺序 & 典型风险 \\
\midrule
所有边代价相同 & BFS & 队列按层 & 边权一变就失效 \\
边权只有 0 和 1 & 0-1 BFS & 0 入队首，1 入队尾 & 不能只用普通 \code{visited} \\
边权非负 & Dijkstra & 小顶堆取最小距离 & 负边会破坏弹出确定性 \\
负边或限制边数 & Bellman-Ford & 按边逐轮松弛 & 原地更新会串阶段 \\
多点对密集查询 & Floyd-Warshall & 枚举中间点 & 只适合较小节点数 \\
\bottomrule
\end{tabular}
\end{center}

这条线不是“越来越高级所以越来越好”，而是适用条件逐渐放宽、代价也逐渐变化。边权全为 1 时，Dijkstra 当然也能跑，但堆是多余的；边权只有 0 和 1 时，0-1 BFS 比堆更贴合结构；存在负边时，Dijkstra 的弹出确定性失效，必须换成按边数逐轮传播的思想。好的题解应该说明为什么选这个算法，而不只是把代码贴出来。

网络延迟时间是 Dijkstra 的典型题。节点是服务器，有向边带传播时间。题目问从起点发信号，到所有节点都收到信号需要多久，也就是起点到所有节点最短路的最大值。暴力方法可以反复松弛所有边，类似 Bellman-Ford，复杂度 \complexity{O(nm)}。由于边权非负，用堆优化 Dijkstra 更合适。

\begin{lstlisting}[language=C++,caption={LeetCode 743 网络延迟时间：堆优化 Dijkstra}]
int network_delay_time(const std::vector<std::vector<int>>& times,
                       int n,
                       int source)
{
    using Edge = std::pair<int, int>;  // neighbor, weight
    std::vector<std::vector<Edge>> graph(n + 1);
    for (const std::vector<int>& edge : times) {
        graph[edge[0]].emplace_back(edge[1], edge[2]);
    }

    constexpr int INF = 1'000'000'000;
    std::vector<int> distance(n + 1, INF);
    distance[source] = 0;

    using State = std::pair<int, int>;  // distance, node
    std::priority_queue<State, std::vector<State>, std::greater<State>> heap;
    heap.emplace(0, source);

    while (!heap.empty()) {
        const auto [dist, node] = heap.top();
        heap.pop();
        if (dist != distance[node]) {
            continue;
        }

        for (const auto [next, weight] : graph[node]) {
            if (dist + weight < distance[next]) {
                distance[next] = dist + weight;
                heap.emplace(distance[next], next);
            }
        }
    }

    int answer = 0;
    for (int node = 1; node <= n; ++node) {
        if (distance[node] == INF) {
            return -1;
        }
        answer = std::max(answer, distance[node]);
    }
    return answer;
}
\end{lstlisting}

这段代码没有使用显式 \code{visited}，而是用 \code{dist != distance[node]} 跳过堆中的过期状态。因为 C++ 标准堆不支持 decrease-key，更新距离时直接把新状态再压入堆，旧状态留在堆里，弹出时识别并跳过。复杂度是 \complexity{O((n+m)\log m)} 或常写为 \complexity{O(m\log n)}。如果存在负边，Dijkstra 不适用，应考虑 Bellman-Ford 或 SPFA 的风险分析。

Dijkstra 的证明可以用“反证 + 非负边”来讲。设当前从堆中弹出的节点是 \code{u}，距离为 \code{dist[u]}，如果未来还有一条更短路径能到达 \code{u}，这条路径必须先从已确定区域走到某个未确定节点 \code{x}，再继续走到 \code{u}。由于边权非负，从起点到 \code{x} 的距离不会大于整条更短路径的距离，因此 \code{x} 应该比 \code{u} 更早被弹出，矛盾。这个证明也解释了为什么负边会破坏 Dijkstra：负边可能让“先到某个未确定节点再绕回来”的总代价突然变小。

在 C++ 实现中还有两个工程细节。第一，距离数组要根据题目范围选择 \code{int} 或 \code{std::int64_t}；边数和边权都大时，\code{int} 容易溢出。第二，堆里的状态排序必须以当前路径代价为第一关键字，不能只按节点编号，也不能把 \code{pair} 写反。写 \code{using State = std::pair<int, int>} 时，要明确注释是 \code{distance, node}，否则读者很容易把 \code{node, distance} 压入堆，代码仍能编译但语义错了。

最小体力消耗路径看似网格题，实质是最短路变体。路径代价不是边权之和，而是路径上最大高度差的最小可能值。状态是格子，边是上下左右相邻格子，走一条边的代价是高度差。到达某个格子的最优体力定义为：从起点到它的路径中，最大边权尽可能小。转移时，新代价是 \code{max(当前路径体力, 新边高度差)}。这个代价仍满足 Dijkstra 的非负和单调扩展性质。

\begin{lstlisting}[language=C++,caption={LeetCode 1631 最小体力路径：以最大边权作为路径代价}]
int minimum_effort_path(const std::vector<std::vector<int>>& heights)
{
    constexpr std::array<std::pair<int, int>, 4> DIRECTIONS{
        std::pair{1, 0}, std::pair{-1, 0},
        std::pair{0, 1}, std::pair{0, -1}
    };
    const int rows = static_cast<int>(heights.size());
    const int cols = static_cast<int>(heights[0].size());
    constexpr int INF = 1'000'000'000;

    std::vector<std::vector<int>> effort(rows, std::vector<int>(cols, INF));
    effort[0][0] = 0;

    using State = std::tuple<int, int, int>;  // effort, row, col
    std::priority_queue<State, std::vector<State>, std::greater<State>> heap;
    heap.emplace(0, 0, 0);

    while (!heap.empty()) {
        const auto [cost, row, col] = heap.top();
        heap.pop();
        if (row == rows - 1 && col == cols - 1) {
            return cost;
        }
        if (cost != effort[row][col]) {
            continue;
        }

        for (const auto [dr, dc] : DIRECTIONS) {
            const int next_row = row + dr;
            const int next_col = col + dc;
            const bool out_of_bounds =
                next_row < 0 || next_row >= rows ||
                next_col < 0 || next_col >= cols;
            if (out_of_bounds) {
                continue;
            }

            const int edge_cost =
                std::abs(heights[row][col] - heights[next_row][next_col]);
            const int next_cost = std::max(cost, edge_cost);
            if (next_cost < effort[next_row][next_col]) {
                effort[next_row][next_col] = next_cost;
                heap.emplace(next_cost, next_row, next_col);
            }
        }
    }

    return 0;
}
\end{lstlisting}

这题也可以二分答案加 BFS：猜一个最大允许体力，看能否只走高度差不超过该值的边到终点。二分答案复杂度是 \complexity{O(mn\log H)}，Dijkstra 是 \complexity{O(mn\log(mn))}。两种方法都合理，面试时可以先给更统一的 Dijkstra，再说明二分可行的原因是“体力上限越大，可达性越强”，存在单调性。

最短路变体的关键是检查“路径代价是否随扩展单调不下降”。普通加法非负边满足，因为加上一条非负边不会让路径变短；最小体力路径也满足，因为 \code{max(old, edge)} 不会小于 \code{old}；最大概率路径如果直接乘概率，路径概率会不增，通常用大顶堆每次取当前概率最大状态，也可以对概率取负对数变成非负加法最短路。若路径代价可能在扩展后变小，就不能直接用 Dijkstra 弹出确定。

\begin{principlebox}{最短路算法选择标准}
先问边权是否统一。所有边代价相同，用 BFS；边权只有 0 和 1，用 0-1 BFS；边权非负且需要最小总代价，用 Dijkstra；边权可能为负或路径边数被限制，用 Bellman-Ford 思想；节点数较小且要回答大量任意两点查询，用 Floyd-Warshall。不要把“图题”直接等同于某个模板，最短路算法的选择来自边权、查询次数、图规模和路径代价的合成方式。
\end{principlebox}

状态图题还包括打开转盘锁、单词接龙、蛇梯棋等。共同点是输入没有直接给边，需要你根据规则生成邻居。生成邻居时要注意两件事：第一，状态编码要能高效去重；第二，访问标记要在入队时设置，避免同一状态被多个父状态重复加入。无权最短路使用 BFS，第一次到达目标即为最短；带权状态图则要换 Dijkstra 或 A*，不能强行 BFS。

\section{图的表示、层次搜索和拓扑顺序}

真正写图论题之前，还要先决定图如何存储。邻接矩阵适合节点数量较小、需要频繁判断任意两点是否有边的场景，空间是 \complexity{O(n^2)}；邻接表适合稀疏图，空间是 \complexity{O(n+m)}；边数组适合 Kruskal、Bellman-Ford 这类按边扫描的算法；网格图通常不需要显式建图，用坐标和方向数组现场生成邻居即可。很多低质量答案只说“建图”，却没有说明为什么这样建。面试中应该先看输入规模：如果 \code{n} 可以到十万，邻接矩阵基本不可能；如果 \code{m} 远小于 \code{n^2}，邻接表更自然；如果题目要求按边权排序，边数组会比邻接表更直接。

邻接表也有工程差异。 \code{vector<vector<pair<int,int>>>} 写法直观，但每个内层 \code{vector} 都是独立对象，极端情况下会产生很多小分配；竞赛中常见的前向星把所有边放进几个连续数组，通过 \code{head} 和 \code{next} 串起来，局部性更好，但代码可读性低。刷题面试通常以正确性和表达清晰为第一优先，选 \code{vector} 邻接表即可；当题目数据极大或需要讲性能时，可以补充连续边数组的思路。这个判断体现的是工程能力：知道什么时候写简单代码，什么时候需要控制内存布局。

\begin{lstlisting}[language=C++,caption={边数组转邻接表：保留输入语义并便于最短路}]
struct WeightedEdge {
    int from;
    int to;
    int weight;
};

std::vector<std::vector<std::pair<int, int>>> build_directed_graph(
    int node_count,
    const std::vector<WeightedEdge>& edges)
{
    std::vector<std::vector<std::pair<int, int>>> graph(node_count);
    for (const WeightedEdge& edge : edges) {
        graph[edge.from].emplace_back(edge.to, edge.weight);
    }
    return graph;
}
\end{lstlisting}

这段代码看似简单，但它体现了一个重要习惯：输入边和算法图结构分开。 \code{WeightedEdge} 保留题目给出的边语义，邻接表服务于后续算法。若后面同时要跑 Kruskal 和 Dijkstra，就不必从邻接表反推出边数组。很多工程 bug 来自过早把所有信息塞进一种结构，后面算法一换就需要反复转换。

BFS 的本质是按距离层扩展。普通队列能保证这一点，因为所有边权都是 1。访问标记必须在入队时设置，而不是出队时设置；如果出队时才标记，一个节点可能被同一层的多个父节点重复入队，轻则浪费时间，重则在记录父节点或计数时产生错误。多源 BFS 只是把多个起点同时放入队列，把它们的初始距离都设为 0。腐烂的橘子、地图分析、离所有建筑最近的空地等题，都可以用这个模型理解。

\begin{lstlisting}[language=C++,caption={多源 BFS：所有起点同时入队}]
int shortest_distance_from_sources(
    const std::vector<std::vector<int>>& grid,
    const std::vector<std::pair<int, int>>& sources)
{
    constexpr std::array<std::pair<int, int>, 4> DIRECTIONS{
        std::pair{1, 0}, std::pair{-1, 0},
        std::pair{0, 1}, std::pair{0, -1}
    };
    const int rows = static_cast<int>(grid.size());
    const int cols = static_cast<int>(grid[0].size());

    std::vector<std::vector<int>> distance(rows, std::vector<int>(cols, -1));
    std::queue<std::pair<int, int>> queue;
    for (const auto [row, col] : sources) {
        distance[row][col] = 0;
        queue.emplace(row, col);
    }

    int farthest = 0;
    while (!queue.empty()) {
        const auto [row, col] = queue.front();
        queue.pop();
        farthest = std::max(farthest, distance[row][col]);

        for (const auto [dr, dc] : DIRECTIONS) {
            const int next_row = row + dr;
            const int next_col = col + dc;
            const bool out_of_bounds =
                next_row < 0 || next_row >= rows ||
                next_col < 0 || next_col >= cols;
            if (out_of_bounds || grid[next_row][next_col] == 1 ||
                distance[next_row][next_col] != -1) {
                continue;
            }
            distance[next_row][next_col] = distance[row][col] + 1;
            queue.emplace(next_row, next_col);
        }
    }

    return farthest;
}
\end{lstlisting}

这段模板里，\code{distance == -1} 同时承担“尚未访问”和“距离未知”的含义。若题目需要允许障碍、空地、起点不同类型，就不要直接修改原数组，单独开 \code{distance} 会更安全。复杂度是 \complexity{O(mn)}，因为每个格子最多入队一次，每条网格边最多检查常数次。面试时要强调：多源 BFS 不是从每个源分别 BFS，而是把所有源放在第 0 层同时扩散；这样第一次到达某个点的源，就是离它最近的一批源之一。

0-1 BFS 是普通 BFS 和 Dijkstra 之间的桥梁。当边权只可能是 0 或 1 时，使用优先队列当然正确，但 \code{deque} 可以把复杂度降到线性。权重为 0 的转移放到队首，权重为 1 的转移放到队尾。这样队列中的状态仍然按照非降距离处理，只是用双端队列代替了堆。最典型的题是网格中每个格子给出推荐方向，顺着方向走代价为 0，改方向代价为 1，求从左上到右下的最小修改次数。

\begin{lstlisting}[language=C++,caption={LeetCode 1368 使网格图至少有一条有效路径的最小代价：0-1 BFS}]
int min_cost_valid_path(const std::vector<std::vector<int>>& grid)
{
    constexpr std::array<std::pair<int, int>, 4> DIRECTIONS{
        std::pair{0, 1}, std::pair{0, -1},
        std::pair{1, 0}, std::pair{-1, 0}
    };
    constexpr int INF = 1'000'000'000;
    const int rows = static_cast<int>(grid.size());
    const int cols = static_cast<int>(grid[0].size());

    std::vector<std::vector<int>> distance(rows, std::vector<int>(cols, INF));
    std::deque<std::pair<int, int>> deque;
    distance[0][0] = 0;
    deque.emplace_front(0, 0);

    while (!deque.empty()) {
        const auto [row, col] = deque.front();
        deque.pop_front();

        for (int dir = 0; dir < static_cast<int>(DIRECTIONS.size()); ++dir) {
            const auto [dr, dc] = DIRECTIONS[dir];
            const int next_row = row + dr;
            const int next_col = col + dc;
            const bool out_of_bounds =
                next_row < 0 || next_row >= rows ||
                next_col < 0 || next_col >= cols;
            if (out_of_bounds) {
                continue;
            }

            const int extra_cost = (grid[row][col] == dir + 1) ? 0 : 1;
            const int next_distance = distance[row][col] + extra_cost;
            if (next_distance >= distance[next_row][next_col]) {
                continue;
            }

            distance[next_row][next_col] = next_distance;
            if (extra_cost == 0) {
                deque.emplace_front(next_row, next_col);
            } else {
                deque.emplace_back(next_row, next_col);
            }
        }
    }

    return distance[rows - 1][cols - 1];
}
\end{lstlisting}

0-1 BFS 的正确性可以这样解释：普通 BFS 每次扩展当前最小层；Dijkstra 每次从堆里取最小距离；当新增边代价只有 0 或 1 时，当前状态的后继要么仍在当前距离层，要么进入下一距离层。放队首和队尾正好维护了这个层次。易错点有三个：第一，不能用普通 \code{queue}，因为 0 代价边应该优先处理；第二，不能只用 \code{visited}，因为某个状态可能先以较大代价到达，后来被较小代价改进；第三，方向编号要和题目约定对齐，示例里题目方向从 1 开始，所以比较时使用 \code{dir + 1}。

拓扑排序处理有向无环图。它回答的不是最短路，而是依赖顺序：如果一件事必须在另一件事之前完成，就可以画一条有向边。课程表就是标准题。暴力想法是不断找没有前置依赖的课程，删除它，再更新其他课程；拓扑排序把这个过程用入度数组和队列高效实现。若最后处理课程数小于总课程数，说明存在环，环内课程互相等待，无法完成。

\begin{lstlisting}[language=C++,caption={LeetCode 207 课程表：Kahn 拓扑排序判断是否有环}]
bool can_finish(int course_count,
                const std::vector<std::vector<int>>& prerequisites)
{
    std::vector<std::vector<int>> graph(course_count);
    std::vector<int> indegree(course_count, 0);

    for (const std::vector<int>& relation : prerequisites) {
        const int course = relation[0];
        const int dependency = relation[1];
        graph[dependency].push_back(course);
        ++indegree[course];
    }

    std::queue<int> ready;
    for (int course = 0; course < course_count; ++course) {
        if (indegree[course] == 0) {
            ready.push(course);
        }
    }

    int processed = 0;
    while (!ready.empty()) {
        const int course = ready.front();
        ready.pop();
        ++processed;

        for (const int next : graph[course]) {
            --indegree[next];
            if (indegree[next] == 0) {
                ready.push(next);
            }
        }
    }

    return processed == course_count;
}
\end{lstlisting}

这题最容易写反边。题目常写成 \code{[a, b]} 表示想学 \code{a} 必须先学 \code{b}，所以边应该是 \code{b -> a}。入度表示还有多少依赖没有满足。队列中放的是当前已经没有前置依赖的课程。正确性来自不变量：队列中的每门课入度为 0，已经可以安排；处理它等价于完成这门课，因此它指向的后继少一个未满足依赖。若图中有环，环上每个点至少依赖环内另一个点，永远不会全部入队。

拓扑排序不只判断能不能完成，还能在 DAG 上做动态规划。只要图无环，就可以按拓扑顺序保证前驱状态已经计算完，再更新后继。比如有向无环图中求从起点到每个点的最长路径，不能用 Dijkstra，因为目标是最长而不是最短；也不能在有环图上直接做，因为正权环会让最长路径没有上界。DAG 的拓扑顺序让“阶段”变得明确。

\begin{lstlisting}[language=C++,caption={DAG 最长路径：拓扑序上的动态规划}]
std::vector<int> longest_path_in_dag(
    int node_count,
    const std::vector<std::tuple<int, int, int>>& edges,
    int source)
{
    std::vector<std::vector<std::pair<int, int>>> graph(node_count);
    std::vector<int> indegree(node_count, 0);
    for (const auto [from, to, weight] : edges) {
        graph[from].emplace_back(to, weight);
        ++indegree[to];
    }

    std::queue<int> ready;
    for (int node = 0; node < node_count; ++node) {
        if (indegree[node] == 0) {
            ready.push(node);
        }
    }

    constexpr int NEG_INF = -1'000'000'000;
    std::vector<int> distance(node_count, NEG_INF);
    distance[source] = 0;

    while (!ready.empty()) {
        const int node = ready.front();
        ready.pop();
        if (distance[node] != NEG_INF) {
            for (const auto [next, weight] : graph[node]) {
                distance[next] = std::max(distance[next], distance[node] + weight);
            }
        }
        for (const auto [next, weight] : graph[node]) {
            (void)weight;
            --indegree[next];
            if (indegree[next] == 0) {
                ready.push(next);
            }
        }
    }

    return distance;
}
\end{lstlisting}

上面代码故意把“松弛后继”和“减少入度”分开写，便于读者看清两个动作。实际工程中可以合并循环，但教材里先强调概念更重要。DAG DP 的状态是 \code{distance[v]}，含义是从 \code{source} 到 \code{v} 的最大路径和；转移来自所有前驱。拓扑序保证当某个节点出队时，它的所有前驱已经处理完，状态不再被未来节点回头修改。这个结构和一维 DP 很像：所谓“阶段顺序”，在图上就是拓扑顺序。

\section{限制边数、全源最短路和生成树}

如果图中存在负边，Dijkstra 的贪心基础会崩塌。因为一个当前距离较大的节点，未来可能通过负边把另一个已经确定的节点拉得更小。Bellman-Ford 的想法更朴素：一条最短路径如果最多包含 \code{k} 条边，那么做 \code{k} 轮按边松弛就能得到它。经典最短路最多需要 \code{n - 1} 条边，因为无负环的最短简单路径不会重复经过节点。力扣中“最多经过 K 站中转”的机票题，本质是限制边数的 Bellman-Ford。

\begin{lstlisting}[language=C++,caption={LeetCode 787 K 站中转内最便宜航班：按边数分层松弛}]
int find_cheapest_price(int n,
                        const std::vector<std::vector<int>>& flights,
                        int source,
                        int target,
                        int max_stops)
{
    constexpr int INF = 1'000'000'000;
    std::vector<int> previous(n, INF);
    previous[source] = 0;

    for (int step = 0; step <= max_stops; ++step) {
        std::vector<int> current = previous;
        for (const std::vector<int>& flight : flights) {
            const int from = flight[0];
            const int to = flight[1];
            const int price = flight[2];
            if (previous[from] == INF) {
                continue;
            }
            current[to] = std::min(current[to], previous[from] + price);
        }
        previous = std::move(current);
    }

    return previous[target] == INF ? -1 : previous[target];
}
\end{lstlisting}

这里必须使用 \code{previous} 和 \code{current} 两层数组，不能在同一个数组里原地更新。原因是第 \code{step} 轮只允许使用最多 \code{step + 1} 条边，如果原地更新，一轮内刚更新出来的状态又被后面的边继续使用，就会偷偷使用更多边。这个问题和 0-1 背包的倒序遍历本质类似：转移要读取上一阶段，不能让当前阶段污染自己。复杂度是 \complexity{O((K+1)m)}，空间 \complexity{O(n)}。

Bellman-Ford 的完整版本通常做 \code{n - 1} 轮松弛。第 \code{r} 轮结束后，所有使用不超过 \code{r} 条边的最短路径都已经被考虑。为什么是 \code{n - 1}？如果没有负环，一条最短简单路径不会重复经过节点，最多经过 \code{n - 1} 条边。做完 \code{n - 1} 轮后，如果还能继续松弛，说明存在从起点可达的负权环，最短路没有良定义。力扣面试里很少要求完整负环检测，但理解它能帮助你解释为什么 K 站中转题不能用原地 Dijkstra 式确定。

\begin{lstlisting}[language=C++,caption={Bellman-Ford：完整按边松弛和负环检测}]
bool bellman_ford(int node_count,
                  const std::vector<std::tuple<int, int, int>>& edges,
                  int source,
                  std::vector<std::int64_t>& distance)
{
    constexpr std::int64_t INF = 4'000'000'000'000'000'000;
    distance.assign(node_count, INF);
    distance[source] = 0;

    for (int round = 1; round <= node_count - 1; ++round) {
        bool changed = false;
        for (const auto [from, to, weight] : edges) {
            if (distance[from] == INF) {
                continue;
            }
            if (distance[from] + weight < distance[to]) {
                distance[to] = distance[from] + weight;
                changed = true;
            }
        }
        if (!changed) {
            break;
        }
    }

    for (const auto [from, to, weight] : edges) {
        if (distance[from] != INF && distance[from] + weight < distance[to]) {
            return false;
        }
    }
    return true;
}
\end{lstlisting}

这段代码返回 \code{false} 表示存在可达负环。这里使用 \code{std::int64_t} 和很大的 \code{INF}，是为了避免边权累加溢出。与 Dijkstra 相比，Bellman-Ford 不需要邻接表也能工作，因为它每一轮都扫描边数组；这也是为什么“边数组”是图表示的一种独立选择。若题目只限制最多使用 K 条边，就不需要做 \code{n - 1} 轮，而是做 \code{K + 1} 轮，并且为了不串阶段要使用两层数组。

Floyd-Warshall 解决的是所有点对之间的最短路。它的状态可以写成 \code{dist[i][j]}：当前允许使用编号不超过 \code{k} 的中间点时，从 \code{i} 到 \code{j} 的最短距离。每加入一个新的中间点 \code{k}，就尝试把路径拆成 \code{i -> k} 和 \code{k -> j} 两段。复杂度 \complexity{O(n^3)}，空间 \complexity{O(n^2)}，所以它适合节点数较小、查询很多的场景。LeetCode 1334 阈值距离内邻居最少的城市，如果 \code{n} 只有一百量级，Floyd 很直接；如果图很稀疏且节点很多，就更适合从每个点跑 Dijkstra 或只对查询相关点跑最短路。

\begin{lstlisting}[language=C++,caption={LeetCode 1334：Floyd 预处理所有点对距离}]
int find_the_city(int n,
                  const std::vector<std::vector<int>>& edges,
                  int distance_threshold)
{
    constexpr int INF = 1'000'000'000;
    std::vector<std::vector<int>> distance(n, std::vector<int>(n, INF));
    for (int i = 0; i < n; ++i) {
        distance[i][i] = 0;
    }
    for (const std::vector<int>& edge : edges) {
        const int from = edge[0];
        const int to = edge[1];
        const int weight = edge[2];
        distance[from][to] = std::min(distance[from][to], weight);
        distance[to][from] = std::min(distance[to][from], weight);
    }

    for (int mid = 0; mid < n; ++mid) {
        for (int from = 0; from < n; ++from) {
            if (distance[from][mid] == INF) {
                continue;
            }
            for (int to = 0; to < n; ++to) {
                if (distance[mid][to] == INF) {
                    continue;
                }
                distance[from][to] = std::min(
                    distance[from][to],
                    distance[from][mid] + distance[mid][to]);
            }
        }
    }

    int best_city = -1;
    int best_count = INF;
    for (int city = 0; city < n; ++city) {
        int reachable = 0;
        for (int other = 0; other < n; ++other) {
            if (other != city && distance[city][other] <= distance_threshold) {
                ++reachable;
            }
        }
        if (reachable <= best_count) {
            best_count = reachable;
            best_city = city;
        }
    }
    return best_city;
}
\end{lstlisting}

这题的并列规则是选择编号最大的城市，所以更新条件使用 \code{reachable <= best_count}，让后面的较大编号覆盖前面的同数量答案。Floyd 的三层循环顺序不能随便换：最外层必须是中间点 \code{mid}，它表示“当前允许使用哪些中间点”的阶段。如果把 \code{from} 放到最外层，就会让状态含义变乱，可能在同一阶段使用了还不该使用的中间点。

最小生成树回答的是“把所有点连起来的最低总成本”，它不是最短路。最短路关心从某个起点到其他点的路径长度；最小生成树关心选出一组边，让全图连通且总权重最小。Kruskal 的策略是按边权从小到大扫描，能连接两个不同连通块就选，否则跳过。并查集用来判断一条边是否会成环。这个贪心正确性来自割性质：任意一个割上，权重最小的跨割边一定存在某棵最小生成树包含它。

\begin{lstlisting}[language=C++,caption={Kruskal：按边权连接不同连通块}]
struct UndirectedEdge {
    int u;
    int v;
    int weight;
};

int minimum_spanning_tree_cost(int node_count,
                               std::vector<UndirectedEdge> edges)
{
    std::sort(edges.begin(), edges.end(),
              [](const UndirectedEdge& lhs, const UndirectedEdge& rhs) {
                  return lhs.weight < rhs.weight;
              });

    DisjointSetUnion dsu(node_count);
    int total_cost = 0;
    int used_edges = 0;

    for (const UndirectedEdge& edge : edges) {
        if (!dsu.unite(edge.u, edge.v)) {
            continue;
        }
        total_cost += edge.weight;
        ++used_edges;
        if (used_edges == node_count - 1) {
            return total_cost;
        }
    }

    return -1;
}
\end{lstlisting}

面试中区分 Kruskal 和 Prim 的方式很简单：Kruskal 站在“边”的视角，排序后用并查集选边；Prim 站在“点”的视角，从一个连通块开始，每次选连接外部点的最小边。边比较少、已经容易生成边数组时，Kruskal 很顺；完全图且边很多时，Prim 可以避免显式存储所有边。连接所有点的最小费用中，点与点之间距离是曼哈顿距离，朴素建完全图有 \complexity{O(n^2)} 条边，\code{n} 不大时可以接受；如果规模上去，就要研究几何图优化，而不是盲目套模板。

Kruskal 的正确性常用割性质说明。把当前已经选出的边看成若干连通块，任意一个连通块和外部之间形成一个割。跨过这个割的最小边如果连接两个不同连通块，选择它不会破坏最优性，因为任何生成树都必须用至少一条边跨过这个割，而更贵的跨割边可以被这条更便宜的边替换。并查集在这里不是“为了用而用”，它负责回答“这条边是否连接了两个不同连通块”。如果两个端点已经连通，加入它只会形成环，不会让连通性更好。

Prim 的正确性也来自割性质，只是它始终维护一个已经连通的点集。每一步从点集到外部的所有边里选最小边，把一个新点接进来。这个过程和 Dijkstra 看起来都使用堆，但语义完全不同：Dijkstra 堆里是从起点到节点的当前最短路径估计，Prim 堆里是把外部点接入当前生成树的候选边成本。不要因为代码都用了 \code{priority_queue} 就把两个算法混成一个。

\begin{lstlisting}[language=C++,caption={LeetCode 1584 连接所有点的最小费用：Prim 不显式建完全图}]
int min_cost_connect_points(const std::vector<std::vector<int>>& points)
{
    const int n = static_cast<int>(points.size());
    std::vector<int> min_edge(n, std::numeric_limits<int>::max());
    std::vector<char> used(n, false);
    min_edge[0] = 0;

    int total_cost = 0;
    for (int step = 0; step < n; ++step) {
        int node = -1;
        for (int candidate = 0; candidate < n; ++candidate) {
            if (!used[candidate] &&
                (node == -1 || min_edge[candidate] < min_edge[node])) {
                node = candidate;
            }
        }

        used[node] = true;
        total_cost += min_edge[node];

        for (int next = 0; next < n; ++next) {
            if (used[next]) {
                continue;
            }
            const int cost = std::abs(points[node][0] - points[next][0]) +
                std::abs(points[node][1] - points[next][1]);
            min_edge[next] = std::min(min_edge[next], cost);
        }
    }
    return total_cost;
}
\end{lstlisting}

这段 Prim 是 \complexity{O(n^2)}，没有显式存储 \complexity{O(n^2)} 条边，适合点数不大、任意两点边权可以现场计算的完全图。若已经有稀疏边表，Kruskal 可能更自然；若图是稠密图，朴素 Prim 的数组扫描反而比维护大量堆状态更简单稳定。算法选择不是背“哪个更高级”，而是看边是否容易枚举、图是否稠密、是否需要动态生成边权。

并查集还能和二分答案结合。最小体力路径除了 Dijkstra，也可以二分一个体力上限 \code{limit}，只把高度差不超过 \code{limit} 的相邻格子合并，最后看起点和终点是否连通。若某个 \code{limit} 可行，那么更大的上限一定可行；若不可行，更小的上限一定不可行。这就是二分的单调性。这个方法的优点是解释直观，缺点是每次检查都要重新扫描网格。

\begin{lstlisting}[language=C++,caption={LeetCode 1631 最小体力路径：二分答案加并查集}]
bool can_reach_with_effort(const std::vector<std::vector<int>>& heights,
                           int limit)
{
    const int rows = static_cast<int>(heights.size());
    const int cols = static_cast<int>(heights[0].size());
    DisjointSetUnion dsu(rows * cols);

    auto id = [cols](int row, int col) {
        return row * cols + col;
    };

    for (int row = 0; row < rows; ++row) {
        for (int col = 0; col < cols; ++col) {
            if (row + 1 < rows &&
                std::abs(heights[row][col] - heights[row + 1][col]) <= limit) {
                dsu.unite(id(row, col), id(row + 1, col));
            }
            if (col + 1 < cols &&
                std::abs(heights[row][col] - heights[row][col + 1]) <= limit) {
                dsu.unite(id(row, col), id(row, col + 1));
            }
        }
    }

    return dsu.find(0) == dsu.find(rows * cols - 1);
}

int minimum_effort_path_by_binary_search(
    const std::vector<std::vector<int>>& heights)
{
    int low = 0;
    int high = 1'000'000;
    while (low < high) {
        const int mid = low + (high - low) / 2;
        if (can_reach_with_effort(heights, mid)) {
            high = mid;
        } else {
            low = mid + 1;
        }
    }
    return low;
}
\end{lstlisting}

这段代码只合并向下和向右的边，是因为无向网格中每条相邻边检查一次即可。复杂度是 \complexity{O(mn\log H \cdot \alpha(mn))}，其中 \code{H} 是高度差范围。和 Dijkstra 相比，二分并查集依赖答案单调性；Dijkstra 依赖路径代价的单调扩展性质。两种方法各有适用场景。面试时能同时说出这两种建模，说明你不是在背题，而是在识别“最小化最大值”的结构。

\section{强连通、二分图匹配和网络流}

有些图题的答案不在一条路径上，而在图的整体结构上。无向图里，某条边一旦删除就让连通块数量增加，它就是桥；某个点一旦删除就让连通块数量增加，它就是割点。有向图里，若一组点两两可达，它们形成强连通分量；把每个强连通分量缩成一个点后，原图会变成有向无环图。二分图匹配关心的是两侧对象如何一对一配对；网络流则把“容量约束下最多能送多少单位”作为统一模型。这些算法看起来比 BFS、Dijkstra 更高级，但它们的核心仍然是状态和不变量。

先看强连通分量。暴力方法可以从每个点出发做可达性搜索，再判断哪些点互相可达；这会重复扫描大量边。强连通分量算法利用一个结构事实：如果把每个强连通分量缩成一个点，得到的分量图一定是 DAG。Kosaraju 算法分两遍搜索：第一遍在原图上按完成时间得到一个顺序；第二遍在反图上按这个顺序反向扩展，每次扩展得到一个强连通分量。完成时间的作用是把分量图中的“源”或“汇”按合适顺序剥离出来。

\begin{lstlisting}[language=C++,caption={Kosaraju：显式栈求有向图强连通分量}]
std::vector<int> strongly_connected_components(
    const std::vector<std::vector<int>>& graph)
{
    const int n = static_cast<int>(graph.size());
    std::vector<std::vector<int>> reversed(n);
    for (int from = 0; from < n; ++from) {
        for (const int to : graph[from]) {
            reversed[to].push_back(from);
        }
    }

    std::vector<char> visited(n, false);
    std::vector<int> order;
    order.reserve(n);

    for (int start = 0; start < n; ++start) {
        if (visited[start]) {
            continue;
        }
        std::vector<std::pair<int, int>> stack{{start, 0}};
        visited[start] = true;
        while (!stack.empty()) {
            auto& frame = stack.back();
            const int node = frame.first;
            int& next_index = frame.second;
            if (next_index < static_cast<int>(graph[node].size())) {
                const int next = graph[node][next_index++];
                if (!visited[next]) {
                    visited[next] = true;
                    stack.emplace_back(next, 0);
                }
                continue;
            }
            order.push_back(node);
            stack.pop_back();
        }
    }

    std::vector<int> component(n, -1);
    int component_count = 0;
    for (int i = n - 1; i >= 0; --i) {
        const int start = order[i];
        if (component[start] != -1) {
            continue;
        }

        std::vector<int> stack{start};
        component[start] = component_count;
        while (!stack.empty()) {
            const int node = stack.back();
            stack.pop_back();
            for (const int next : reversed[node]) {
                if (component[next] != -1) {
                    continue;
                }
                component[next] = component_count;
                stack.push_back(next);
            }
        }
        ++component_count;
    }

    return component;
}
\end{lstlisting}

这段代码没有使用递归。第一遍显式保存 \code{(node, next_index)}，模拟递归 DFS 的“进入节点、逐条边处理、所有边处理完再记录完成时间”。第二遍在反图上普通栈扩展，把能到达的点标成同一个分量。强连通分量常用于“最少加多少边让图强连通”“有向图中哪些点互相可达”“2-SAT 可满足性”等问题。刷题中如果题目只要求无向连通块，普通 BFS 或并查集就够；只有有向互达关系才需要强连通分量。

桥的核心概念是 \code{low} 值。对无向图做 DFS 时，\code{disc[u]} 是节点第一次被访问的时间，\code{low[u]} 是从 \code{u} 的子树出发，最多走一条返祖边能到达的最早时间。若树边 \code{parent -> child} 满足 \code{low[child] > disc[parent]}，说明 \code{child} 子树无法通过其他边回到 \code{parent} 或更早祖先；删除这条边会把图切开，所以它是桥。这个判断不是模板魔法，而是在问“子树是否还有备用通道回到上层”。

\begin{lstlisting}[language=C++,caption={LeetCode 1192 关键连接：显式栈 Tarjan 桥算法}]
std::vector<std::vector<int>> critical_connections(
    int node_count,
    const std::vector<std::vector<int>>& connections)
{
    struct Edge {
        int to;
        int id;
    };

    std::vector<std::vector<Edge>> graph(node_count);
    for (int id = 0; id < static_cast<int>(connections.size()); ++id) {
        const int u = connections[id][0];
        const int v = connections[id][1];
        graph[u].push_back(Edge{v, id});
        graph[v].push_back(Edge{u, id});
    }

    std::vector<int> discover_time(node_count, -1);
    std::vector<int> low_time(node_count, 0);
    std::vector<int> parent(node_count, -1);
    std::vector<int> parent_edge(node_count, -1);
    std::vector<std::vector<int>> bridges;

    struct Frame {
        int node;
        int next_index;
    };

    int timer = 0;
    for (int start = 0; start < node_count; ++start) {
        if (discover_time[start] != -1) {
            continue;
        }

        discover_time[start] = timer;
        low_time[start] = timer;
        ++timer;
        std::vector<Frame> stack{{start, 0}};

        while (!stack.empty()) {
            Frame& frame = stack.back();
            const int node = frame.node;
            if (frame.next_index < static_cast<int>(graph[node].size())) {
                const Edge edge = graph[node][frame.next_index++];
                if (edge.id == parent_edge[node]) {
                    continue;
                }
                const int next = edge.to;
                if (discover_time[next] == -1) {
                    parent[next] = node;
                    parent_edge[next] = edge.id;
                    discover_time[next] = timer;
                    low_time[next] = timer;
                    ++timer;
                    stack.push_back(Frame{next, 0});
                } else {
                    low_time[node] =
                        std::min(low_time[node], discover_time[next]);
                }
                continue;
            }

            stack.pop_back();
            const int previous = parent[node];
            if (previous == -1) {
                continue;
            }
            low_time[previous] = std::min(low_time[previous], low_time[node]);
            if (low_time[node] > discover_time[previous]) {
                bridges.push_back({previous, node});
            }
        }
    }

    return bridges;
}
\end{lstlisting}

桥算法最容易错在“跳过父边”的细节。上面使用边编号，而不是只判断 \code{next == parent[node]}，是为了正确处理重边：如果两个点之间有两条平行边，沿其中一条进入子节点，另一条仍然是一条返祖边，它会让这条连接不再是桥。割点的思想也用 \code{low} 值，只是判断对象从边变成点：若某个非根节点存在子节点 \code{child} 使得 \code{low[child] >= disc[node]}，删除该点会让子树无法回到更早祖先；根节点则需要至少两个 DFS 子树才是割点。

二分图的第一层问题是染色。若图可以把节点分成左右两侧，并且所有边都跨侧连接，它就是二分图。暴力判断可能会枚举所有划分，复杂度指数级；染色法从任意未染色点出发，把相邻点染成相反颜色，若某条边两端颜色相同，就不是二分图。这个判断的本质是检查图中是否存在奇环。奇环无法二染色，偶环可以。

二分图的第二层问题是最大匹配。匹配要求每个左侧点最多匹配一个右侧点，每个右侧点也最多匹配一个左侧点。贪心地给每个左侧点选择第一个可用右侧点并不可靠，因为早期选择可能挡住后面唯一可行的点。增广路算法的想法是：如果当前有一个未匹配左点，尝试沿着“未匹配边、已匹配边、未匹配边”交替前进，找到一个未匹配右点；一旦找到，就把路径上的匹配状态全部翻转，匹配数增加一。

\begin{lstlisting}[language=C++,caption={LeetCode 1820 邀请的最大数量：BFS 增广路二分图匹配}]
int maximum_invitations(const std::vector<std::vector<int>>& grid)
{
    const int left_count = static_cast<int>(grid.size());
    const int right_count = static_cast<int>(grid[0].size());
    std::vector<int> match_left(left_count, -1);
    std::vector<int> match_right(right_count, -1);

    int answer = 0;
    for (int start = 0; start < left_count; ++start) {
        if (match_left[start] != -1) {
            continue;
        }

        std::queue<int> queue;
        std::vector<char> seen_left(left_count, false);
        std::vector<char> seen_right(right_count, false);
        std::vector<int> parent_right(right_count, -1);

        queue.push(start);
        seen_left[start] = true;
        int free_right = -1;

        while (!queue.empty() && free_right == -1) {
            const int left = queue.front();
            queue.pop();

            for (int right = 0; right < right_count; ++right) {
                if (grid[left][right] == 0 || seen_right[right]) {
                    continue;
                }
                seen_right[right] = true;
                parent_right[right] = left;

                if (match_right[right] == -1) {
                    free_right = right;
                    break;
                }

                const int next_left = match_right[right];
                if (!seen_left[next_left]) {
                    seen_left[next_left] = true;
                    queue.push(next_left);
                }
            }
        }

        if (free_right == -1) {
            continue;
        }

        int right = free_right;
        while (right != -1) {
            const int left = parent_right[right];
            const int previous_right = match_left[left];
            match_left[left] = right;
            match_right[right] = left;
            right = previous_right;
        }
        ++answer;
    }

    return answer;
}
\end{lstlisting}

这段代码的关键是“翻转路径”。假设路径是 \code{L0 - R0 - L1 - R1}，其中 \code{R1} 未匹配，\code{R0} 原来匹配 \code{L1}。翻转后，\code{L1} 改匹配 \code{R1}，\code{L0} 匹配 \code{R0}，匹配总数增加一。这个过程不是局部贪心，而是在为早期选择寻找调整空间。更高效的 Hopcroft-Karp 算法一次 BFS 分层后寻找多条最短增广路，可以把复杂度降到 \complexity{O(E\sqrt{V})}；面试中多数题用上面的增广路思想已经足够，但要能说明为什么简单贪心不可靠。

网络流把二分图匹配继续抽象。给定源点、汇点和每条边容量，问从源点最多能向汇点送多少流量。每次沿一条仍有剩余容量的路径增广，把这条路径能承受的最小剩余容量加到总流量里，同时在反向边上增加同样容量，表示未来可以撤销或调整这次选择。反向边是网络流最重要的思想之一：它让算法不是一次性贪心，而是允许后续路径修正已有流量分配。

\begin{lstlisting}[language=C++,caption={Edmonds-Karp：BFS 增广路最大流入口实现}]
class FlowNetwork {
public:
    struct FlowEdge {
        int to;
        int reverse_index;
        int capacity;
    };

    explicit FlowNetwork(int node_count) : graph_(node_count) {}

    void add_edge(int from, int to, int capacity)
    {
        FlowEdge forward{
            to, static_cast<int>(this->graph_[to].size()), capacity};
        FlowEdge backward{
            from, static_cast<int>(this->graph_[from].size()), 0};
        this->graph_[from].push_back(forward);
        this->graph_[to].push_back(backward);
    }

    int max_flow(int source, int sink)
    {
        constexpr int INF = 1'000'000'000;
        int total_flow = 0;

        while (true) {
            std::vector<int> parent_node(this->graph_.size(), -1);
            std::vector<int> parent_edge(this->graph_.size(), -1);
            std::queue<int> queue;
            queue.push(source);
            parent_node[source] = source;

            while (!queue.empty() && parent_node[sink] == -1) {
                const int node = queue.front();
                queue.pop();
                for (int edge_index = 0;
                     edge_index < static_cast<int>(this->graph_[node].size());
                     ++edge_index) {
                    const FlowEdge& edge = this->graph_[node][edge_index];
                    if (edge.capacity <= 0 || parent_node[edge.to] != -1) {
                        continue;
                    }
                    parent_node[edge.to] = node;
                    parent_edge[edge.to] = edge_index;
                    queue.push(edge.to);
                    if (edge.to == sink) {
                        break;
                    }
                }
            }

            if (parent_node[sink] == -1) {
                break;
            }

            int pushed = INF;
            for (int node = sink; node != source; node = parent_node[node]) {
                const FlowEdge& edge =
                    this->graph_[parent_node[node]][parent_edge[node]];
                pushed = std::min(pushed, edge.capacity);
            }

            for (int node = sink; node != source; node = parent_node[node]) {
                FlowEdge& edge =
                    this->graph_[parent_node[node]][parent_edge[node]];
                edge.capacity -= pushed;
                this->graph_[node][edge.reverse_index].capacity += pushed;
            }
            total_flow += pushed;
        }

        return total_flow;
    }

private:
    std::vector<std::vector<FlowEdge>> graph_;
};
\end{lstlisting}

Edmonds-Karp 每次用 BFS 找最短增广路，复杂度是 \complexity{O(VE^2)}，不是工业级最大流的终点。Dinic 会先用 BFS 建层次图，再在层次图上推阻塞流，通常更快；费用流还会给边加费用，解决“最大数量中成本最小”的问题。教材这里先给 Edmonds-Karp，是为了把残量网络、反向边、增广路径讲清楚。二分图最大匹配可以看成一个特殊网络流：源点连所有左侧点，容量为 1；左侧按可匹配关系连右侧，容量为 1；右侧连汇点，容量为 1。最大流量就是最大匹配数。

\begin{principlebox}{结构性图算法的选型}
只问无向连通性，用并查集或 BFS；问删除边或点是否破坏连通性，用 \code{disc/low} 和桥、割点；问有向互相可达，用强连通分量；问二分图能否分成两侧，用染色；问两侧对象最多配多少对，用二分图匹配；问容量约束下能通过多少流量，用网络流。结构性图题不要急着套最短路，因为答案往往不是路径长度，而是分量、割、匹配或流量。
\end{principlebox}

\section{状态图、双向搜索和支配剪枝}

账户合并是并查集的另一个高频场景。题目给出若干账户，每个账户包含姓名和邮箱；只要两个账户共享邮箱，就属于同一个人。暴力方法是两两比较账户邮箱集合，复杂度可能达到 \complexity{O(n^2s)}，其中 \code{s} 是邮箱数量。更好的模型是：邮箱是节点，同一个账户里的所有邮箱属于同一集合。扫描账户时，把第一个邮箱和后续邮箱合并；最后按根收集邮箱，再排序输出。

\begin{lstlisting}[language=C++,caption={LeetCode 721 账户合并：把邮箱映射为并查集节点}]
std::vector<std::vector<std::string>> accounts_merge(
    const std::vector<std::vector<std::string>>& accounts)
{
    std::unordered_map<std::string, int> email_id;
    std::unordered_map<std::string, std::string> email_name;

    for (const std::vector<std::string>& account : accounts) {
        const std::string& name = account[0];
        for (int i = 1; i < static_cast<int>(account.size()); ++i) {
            const std::string& email = account[i];
            if (!email_id.contains(email)) {
                const int id = static_cast<int>(email_id.size());
                email_id.emplace(email, id);
            }
            email_name[email] = name;
        }
    }

    DisjointSetUnion dsu(static_cast<int>(email_id.size()));
    for (const std::vector<std::string>& account : accounts) {
        const int first_id = email_id[account[1]];
        for (int i = 2; i < static_cast<int>(account.size()); ++i) {
            dsu.unite(first_id, email_id[account[i]]);
        }
    }

    std::unordered_map<int, std::vector<std::string>> groups;
    for (const auto& [email, id] : email_id) {
        groups[dsu.find(id)].push_back(email);
    }

    std::vector<std::vector<std::string>> answer;
    answer.reserve(groups.size());
    for (auto& [root, emails] : groups) {
        (void)root;
        std::sort(emails.begin(), emails.end());
        std::vector<std::string> merged;
        merged.reserve(emails.size() + 1);
        merged.push_back(email_name[emails.front()]);
        merged.insert(merged.end(), emails.begin(), emails.end());
        answer.push_back(std::move(merged));
    }
    return answer;
}
\end{lstlisting}

这里的关键不是并查集代码，而是节点选择。若把账户编号作为节点，也能合并共享邮箱的账户，但需要用邮箱反查已经出现过的账户编号；若把邮箱作为节点，则最后天然得到邮箱分组。由于输出要求邮箱有序，每个组还要排序。姓名通常不用于判断合并，只用于输出；如果真实业务中同一邮箱对应多个姓名，就要定义冲突规则，不能默认题目条件成立。刷题题解可以借题目保证，工程代码必须把数据一致性说清楚。

状态图建模有一个核心原则：如果两个状态未来可选动作完全一样，并且目标函数只取决于未来，那么它们才能合并。账户合并中，同一个邮箱节点只代表身份关系，路径长短不重要，所以可以合并到同一集合；最短路径消除障碍物中，同一个格子但剩余消除次数不同，未来可穿越障碍能力不同，所以不能合并成一个 \code{visited[row][col]}。很多高级图题不是难在 BFS 代码，而是难在判断哪些历史信息必须成为状态维度。

单词接龙是状态图建模的代表。每个单词是节点，若两个单词只差一个字符，就有一条无权边。暴力建图需要两两比较单词，复杂度高。更好的做法是用通配模式建立中间索引，例如 \code{hot} 可以生成 \code{*ot}、\code{h*t}、\code{ho*}。所有拥有同一通配模式的单词互为相邻候选。搜索时从当前单词生成模式，再找到候选单词。由于边权相同，用 BFS，第一次到达目标就是最短转换长度。

\begin{lstlisting}[language=C++,caption={LeetCode 127 单词接龙：通配模式建索引后 BFS}]
int ladder_length(const std::string& begin_word,
                  const std::string& end_word,
                  const std::vector<std::string>& word_list)
{
    std::unordered_set<std::string> dictionary(word_list.begin(), word_list.end());
    if (!dictionary.contains(end_word)) {
        return 0;
    }
    dictionary.insert(begin_word);

    std::unordered_map<std::string, std::vector<std::string>> buckets;
    for (const std::string& word : dictionary) {
        for (int i = 0; i < static_cast<int>(word.size()); ++i) {
            std::string pattern = word;
            pattern[i] = '*';
            buckets[pattern].push_back(word);
        }
    }

    std::queue<std::string> queue;
    std::unordered_map<std::string, int> distance;
    queue.push(begin_word);
    distance[begin_word] = 1;

    while (!queue.empty()) {
        const std::string word = queue.front();
        queue.pop();
        if (word == end_word) {
            return distance[word];
        }

        for (int i = 0; i < static_cast<int>(word.size()); ++i) {
            std::string pattern = word;
            pattern[i] = '*';
            for (const std::string& next : buckets[pattern]) {
                if (distance.contains(next)) {
                    continue;
                }
                distance[next] = distance[word] + 1;
                queue.push(next);
            }
        }
    }

    return 0;
}
\end{lstlisting}

这段代码为了表达清楚保留了字符串拷贝。面试里可以进一步优化：每个模式桶被访问后清空，避免不同路径反复扫描同一桶；也可以使用双向 BFS，从起点集合和终点集合中较小的一侧扩展，降低搜索宽度。暴力为什么慢？因为单词数量为 \code{n}、长度为 \code{L} 时，两两比较是 \complexity{O(n^2L)}；通配索引构建约为 \complexity{O(nL^2)}，搜索时主要消耗在模式生成和桶扫描。若使用 \code{string_view} 或预编码单词，工程上还能减少临时字符串分配，但面试中要先保证模型清楚。

通配模式索引的本质是把“比较两个单词是否只差一位”改写成“两个单词是否落入同一个桶”。暴力比较每对单词时，同一个前缀和后缀关系会被重复检查；桶索引把这些关系预处理出来，搜索时只查与当前词相关的桶。这个思想可以迁移到很多题：如果判断两个对象是否相邻的成本高，就尝试建立中间 key，把可能相邻的对象提前聚类。缺点是桶可能很大，某些模式会连接大量单词，所以访问后清空桶是一种重要优化，避免同一个桶被不同路径重复扫描。

双向 BFS 的优化依据是搜索树宽度。单向 BFS 深度为 \code{d}、平均分支因子为 \code{b} 时，可能访问约 \code{b^d} 个状态；双向 BFS 从两端各扩展约 \code{d/2} 层，理想情况下访问量接近 \code{2b^{d/2}}。这不是改变最坏复杂度的魔法，而是利用了许多状态图的指数扩展特征。写双向 BFS 时，永远扩展当前较小的一侧，并在生成新状态时检查它是否已经出现在另一侧。如果相遇，当前步数就是答案。

\begin{lstlisting}[language=C++,caption={LeetCode 752 打开转盘锁：双向 BFS 的状态集合扩展}]
int open_lock(const std::vector<std::string>& deadends,
              const std::string& target)
{
    const std::unordered_set<std::string> dead(deadends.begin(), deadends.end());
    const std::string start = "0000";
    if (dead.contains(start)) {
        return -1;
    }
    if (target == start) {
        return 0;
    }

    auto neighbors = [](const std::string& state) {
        std::vector<std::string> result;
        result.reserve(8);
        for (int i = 0; i < static_cast<int>(state.size()); ++i) {
            for (const int delta : {-1, 1}) {
                std::string next = state;
                const int digit = state[i] - '0';
                next[i] = static_cast<char>('0' + (digit + delta + 10) % 10);
                result.push_back(std::move(next));
            }
        }
        return result;
    };

    std::unordered_set<std::string> front{start};
    std::unordered_set<std::string> back{target};
    std::unordered_set<std::string> visited{start, target};
    int steps = 0;

    while (!front.empty() && !back.empty()) {
        if (front.size() > back.size()) {
            std::swap(front, back);
        }
        std::unordered_set<std::string> next_front;
        ++steps;

        for (const std::string& state : front) {
            for (const std::string& next : neighbors(state)) {
                if (dead.contains(next) || visited.contains(next)) {
                    continue;
                }
                if (back.contains(next)) {
                    return steps;
                }
                visited.insert(next);
                next_front.insert(next);
            }
        }
        front = std::move(next_front);
    }

    return -1;
}
\end{lstlisting}

转盘锁的状态空间最多一万个，因此普通 BFS 也足够；这里展示双向 BFS，是为了训练状态集合扩展的写法。注意 \code{visited} 同时包含两端访问过的状态，避免重复回头。若先检查 \code{visited} 再检查 \code{back}，可能会错过相遇状态，所以上面代码对新状态的处理顺序是：死点跳过，已访问跳过，相遇返回，再加入下一层。也可以把相遇检查放在已访问之前，并单独维护两侧访问集；关键是语义一致。

双向 BFS 也有边界。第一，它要求你知道终点集合，并且反向扩展的规则容易生成；如果只有一个起点、终点是“任意满足条件的状态”，双向搜索不一定自然。第二，边最好是无向或可逆；如果边是有向的，反向扩展必须使用反向边，否则相遇判断不成立。第三，双向 BFS 通常只求最短步数，不直接给出完整路径；若要还原路径，需要分别记录两侧父指针，并在相遇点拼接。把这些限制讲清楚，比简单说“双向 BFS 更快”更有价值。

图题中还有一类“边不是原题里的动作，而是你定义出来的关系”。比如岛屿数量把相邻陆地视为同一连通块；方程除法把变量视为节点、比例视为带权有向边；除法求值时，\code{a / b = 2} 可以建立 \code{a -> b} 权重 2 和 \code{b -> a} 权重 \code{1/2}，查询 \code{x / y} 就是从 \code{x} 到 \code{y} 的路径乘积。这个模型不适合普通最短路，因为路径权重不是相加，而是相乘；但 DFS/BFS 沿路径累计乘积即可。若查询很多，还可以用带权并查集维护每个节点到根的比例。

带权并查集比普通并查集难，因为集合关系不仅是“是否连通”，还要维护相对权值。以方程除法为例，\code{weight[x]} 可以表示 \code{x / parent[x]}。路径压缩时，需要同时更新 \code{x / root}。合并 \code{x} 和 \code{y} 且已知 \code{x / y = value} 时，若令 \code{root_x} 挂到 \code{root_y}，就要计算 \code{root_x / root_y} 的值。这个推导容易错，因此面试中如果时间紧，先给 BFS 图解法更稳；只有查询量大或题目明确动态合并时，再写带权并查集。

网格图的另一个高频变化是“状态不只包含位置”。最短路径消除障碍物中，状态应为 \code{(row, col, remaining)}，表示走到某个格子时还剩多少次消除机会。若只用 \code{visited[row][col]}，会错误地剪掉更优状态：同样到达一个格子，剩余消除次数越多，后续能力越强。正确的访问标记可以是三维布尔数组，也可以记录到达每个格子时见过的最大剩余次数；当新状态剩余次数不超过历史最好时，才可以剪掉。

\begin{lstlisting}[language=C++,caption={LeetCode 1293 网格中的最短路径：位置加资源才是完整状态}]
int shortest_path_with_obstacle_elimination(
    const std::vector<std::vector<int>>& grid,
    int eliminations)
{
    constexpr std::array<std::pair<int, int>, 4> DIRECTIONS{
        std::pair{1, 0}, std::pair{-1, 0},
        std::pair{0, 1}, std::pair{0, -1}
    };
    const int rows = static_cast<int>(grid.size());
    const int cols = static_cast<int>(grid[0].size());

    struct State {
        int row;
        int col;
        int remaining;
        int steps;
    };

    std::vector<std::vector<int>> best_remaining(
        rows, std::vector<int>(cols, -1));
    std::queue<State> queue;
    queue.push(State{0, 0, eliminations, 0});
    best_remaining[0][0] = eliminations;

    while (!queue.empty()) {
        const State state = queue.front();
        queue.pop();
        if (state.row == rows - 1 && state.col == cols - 1) {
            return state.steps;
        }

        for (const auto [dr, dc] : DIRECTIONS) {
            const int next_row = state.row + dr;
            const int next_col = state.col + dc;
            const bool out_of_bounds =
                next_row < 0 || next_row >= rows ||
                next_col < 0 || next_col >= cols;
            if (out_of_bounds) {
                continue;
            }
            const int next_remaining =
                state.remaining - grid[next_row][next_col];
            if (next_remaining < 0 ||
                next_remaining <= best_remaining[next_row][next_col]) {
                continue;
            }
            best_remaining[next_row][next_col] = next_remaining;
            queue.push(State{next_row, next_col, next_remaining, state.steps + 1});
        }
    }

    return -1;
}
\end{lstlisting}

这个题非常适合训练“状态支配”思想。到达同一格子时，步数由 BFS 层保证是非降的；如果某个新状态剩余消除次数更少或相等，那么它在未来不会比历史状态更好，可以剪掉。若新状态剩余次数更多，即使步数不更短，也可能穿过更多障碍，不能剪。状态支配比简单 \code{visited} 更细，是很多高阶 BFS 的核心。

支配关系可以用一句话判断：在同一个外部位置上，如果状态 A 的已付成本不高于状态 B，并且 A 的剩余资源、未来选择集合不弱于 B，那么 B 可以删除。障碍消除中，“同格子、更短或同样短、剩余次数更多”支配“同格子、更长且剩余次数更少”；带钥匙迷宫中，钥匙集合大的状态通常支配钥匙集合小的状态，但只有在位置相同且步数不更差时才成立；带时间窗口的图中，早到不一定支配晚到，因为某些门可能要等到特定时间才开。支配剪枝不是随便少存维度，而是要证明未来选择集合真的包含了另一个状态。

\begin{principlebox}{状态图剪枝的判断顺序}
先写完整状态，再尝试压缩。能压缩的前提是被删掉的信息不会影响未来动作、路径代价或答案约束。若要用支配剪枝，必须说明两个状态在同一比较基准上：位置相同、已经付出的成本可比较、剩余资源可比较、未来动作集合存在包含关系。说不出包含关系时，宁可多存一维状态，也不要把不同未来能力错误合并。
\end{principlebox}

\begin{deepdive}
从源码和工程角度看，图算法的性能经常被内存访问模式决定。邻接表遍历时，外层 \code{vector} 的每个桶指向独立存储；节点度数分布很不均匀时，热点节点的邻接边会形成大块连续扫描，而大量低度节点会带来分散访问。前向星或 CSR 风格结构把所有边压进连续数组，遍历时更接近顺序读。Linux 内核大量使用侵入式链表、红黑树、基数树或 XArray 这类结构，是因为内核对象本身已经存在，数据结构常常不拥有对象，只负责把对象链接进某种索引。刷题不需要照搬内核写法，但要理解背后的原则：结构选择服务于访问模式、生命周期和更新频率。
\end{deepdive}

\begin{labbox}
实验建议：选三道题做横向对比。第一道用 Dijkstra 和二分并查集分别解决最小体力路径，记录两种建模的状态、转移和复杂度。第二道用单向 BFS 与双向 BFS 解决打开转盘锁，统计访问状态数量。第三道把课程表从“能否完成”改成“给出一种学习顺序”，再改成“每学期最多学 k 门课时最少需要几学期”，观察拓扑排序、BFS 和状态压缩之间的边界。实验报告不要求长，但必须写清楚为什么一种模型能替代另一种模型，什么时候不能替代。
\end{labbox}

\begin{principlebox}{图论进阶选型}
只问是否连通或连通块数量，优先考虑并查集或 BFS/DFS；问无权最短步数，优先 BFS；问非负权最短距离，优先 Dijkstra；问能否在某个阈值下连通，可以考虑二分答案加 BFS/并查集；问所有边反复松弛且可能有负边，才考虑 Bellman-Ford。先选模型，再写代码。
\end{principlebox}
